%!TEX root=../../main.tex

\section{Equivalence of ReLU and Gated ReLU: Proofs}\label{app:relu-grelu-equivalence}

First we give a simple lemma that will be useful when characterizing the span of \( \calK_i - \calK_i \).

\begin{lemma}\label{lemma:affine-characterization}
    The cone \( \calK_i \) has a non-empty interior if and only if \( \calK_i - \calK_i = \R^d \).
\end{lemma}
\begin{proof}
    Let \( x \in \text{aff} (\calK_i) \).
    Then \( x = \sum_{j=1}^m \alpha_j y_j \), \( y_j \in \calK_i \).
    Let \( j \in [m] \).
    If \( \alpha_j \geq 0 \), then \( \alpha_j y_j \in \calK_i \) since \( \calK_i \) is a cone.
    Otherwise, \( \alpha_j y_j \in - \calK_i \).
    Either way, we have \( \alpha_j y_j \in \calK_i - \calK_i \).

    Observe that \( \calK_i - \calK_i \) is a convex cone since \( \calK_i \) is a convex cone.
    Thus, \( \alpha_1 y_1 + \alpha_2 y_2 \in \calK_i - \calK_i \).
    Induction on \( j \in [m] \) now implies \( x \in \calK_i - \calK_i \) and thus \( \text{aff} (\calK_i) \subseteq \calK_i - \calK_i \).
    Now suppose \( x \in \calK_i - \calK_i \) so that \( x = y_1 - y_2 \), where \( y_1, y_2 \in \calK_i \).
    It is trivial to deduce \( x \in \text{aff} (\calK_i) \);
    we conclude that \( \text{aff} (\calK_i) = \calK_i - \calK_i \).

    Since \( 0 \in \text{aff} (\calK_i) \), this set is a linear subspace of \( \R^d \).
    If \( \calK_i \) has an interior point, then \( \text{aff} (\calK_i) = \R^d \) and we must have \( \text{aff} (\calK_i) = \calK_i - \calK_i = \R^d \).
    On the other hand, if \( \calK_i \) does not have an interior point, then \( \text{aff} (\calK_i) \subset \R^d \) and \( \text{aff} (\calK_i) = \calK_i - \calK_i \subset \R^d \) must hold;
    we have shown the reverse implication by the contrapositive.
\end{proof}

Now we show that \( \calK_i \) has an interior point when \( X \) is full row-rank.
The proof proceeds by studying a relative interior point of \( \calK_i \).

\nonSingularCones*
\begin{proof}
    Let \( \bar w \in \text{relint} (\calK_i) \), which exists since the relative interior of a non-empty convex set is non-empty \citep{bertsekas2009convex}.
    Assume that the inequality,
    \[
        (2 D_i - I) X \bar w \succeq 0,
    \]
    is tight for at least one index \( j \in [n] \);
    let \( X_{\calI} \) be the submatrix of \( X \) formed by the rows of \( X \) for which the inequality
    is tight.
    Define \( \tilde D =  (2 D_i - I) \).
    Since \( X \) is full row-rank, the rows of \( X_{\calI} \) are linearly independent.
    Let \( x_k \) be an arbitrary row of \( X_{\calI} \) (noting \( x_k \neq 0 \) by linear independence) and define \( z_k \) to be the component of \( x_k \) which is orthogonal to the remaining rows of \( X_{\calI} \).
    Clearly such a vector exists since the rows of \( X_{\calI} \) are linearly independent.
    Define \( w' = \bar w + [\tilde D_i]_{kk} z_k \) to obtain
    \[
        [\tilde D_i]_{kk} x_k^\top w' = [\tilde D_i]_{kk} x_k^\top \bar w + \norm{z_k}_2^2 = \norm{z_k}_2^2 > 0,
    \]
    and, for \( j \neq k \),
    \[
        [\tilde D_i]_{jj} x_j^\top w' = [\tilde D_i]_{jj} x_{j}^\top \bar w + [\tilde D_i]_{jj} [D_i]_{kk} x_{j}^\top z_k \geq 0,
    \]
    since \( x_j \) and \( z_k \) are orthogonal.
    This contradicts \( \bar w \in \text{relint} (\calK_i) \) and we deduce that \( (2 D_i - I) X \bar w \succ 0 \).
    \cref{lemma:affine-characterization} now implies that \( \calK_i - \calK_i = \R^d \).

    Let \( u^* \) be an optimal solution to the C-GReLU problem with \( \tilde \calD = \calD_X \).
    Since the Minkowski difference \( \calK_i - \calK_i \) spans \( \R^d \) for every \( D_i \in \calD_X \), we can find \( v_i, u_i \)  such that \( u_i^* = v_i - u_i \).
    Moreover, we can always reparameterize the optimal solution to the C-ReLU problem as \( u_i = v_i^* - u_i^* \).
    A simple reduction argument now shows the two problems are equivalent.
    Applying theorems~\ref{thm:duality-free} and~\ref{thm:unconstrained-equivalence} extends the equivalence to NC-ReLU and NC-GReLU.
\end{proof}

The main difficulty extending \cref{prop:non-singular-cones} to general, full-rank \( X \) is showing that none of the cone-constraints are tight at \( \bar w \).
Unfortunately, the following shows that these difficulties cannot be resolved.

\begin{proposition}\label{prop:col-rank-counter-example}
    There exists a full-rank data matrix \( X \) and activation pattern \( D_i \in \calD_X \) such that \( \calK_i \) is contained in a linear subspace of \( \R^d \).
\end{proposition}
\begin{proof}
    Let \( d = 3 \), \( n = 4 \) and take
    \begin{equation*}
        X =
        \begin{bmatrix}
            1  & 0  & 0  \\
            0  & 1  & 0  \\
            -1 & -1 & 0  \\
            0  & 0  & 1.
        \end{bmatrix}
    \end{equation*}
    It is easy to see that \( X \) is full-rank, although it does not have full row-rank since \( x_1, x_2, x_3 \) are collinear.
    The cone \( \calK_i = \cbr{ w : X w \succeq 0 } \), which corresponds to positive activations for each example, has the following alternative representation:
    \[
        \calK_i = \cbr{\alpha * e_3 : \alpha \geq 0 }.
    \]
    Clearly \( \calK_i \) is contained in a subspace of dimension one.
    Thus, we cannot hope for \( \calK_i \) to have full affine dimension in this more general setting.
\end{proof}

\subsection{Singular Cones are Contained in Non-Singular Cones}\label{app:singular-cones}

Considering the counter-example in \cref{prop:col-rank-counter-example}, we find the ``bad'' \( \calK_i \) is contained within the subspace \( \calS \) spanned by \( e_3 \).
By construction, every \( w \in \calK_i \) is orthogonal to \( x_1, x_2, \) and \( x_3 \), meaning these examples don't contribute to the constraints on \( \calK_i \) once it is restricted to \( \calS \).
Intuitively, changing the activation associated with \( x_1, x_2 \), or \( x_3 \) can only lead to cones which contain \( \calK_i \).
For example, consider \( \calK_i' = \cbr{w : \abr{x_3, w} \leq 0}, X_{-3} w \succeq 0 \), which is equal to the non-negative orthant, \( \R^3_+ \).
We immediately observe \( \calK_i \subset \calK_i' \) and we may replace the degenerate cone with the alternative, full-dimensional \( \calK_2 \).
The rest of this section formalizes these observations.

\begin{definition}\label{def:minimal-constraints}
    Let \( \tilde X \in \R^{m \times d} \) and consider a cone
    \( \calK = \cbr{ w: \tilde X w \succeq 0} \)
    such that \( \text{aff} (\calK) = \calS \subset \R^d \).
    We call an index set \( \calI \subseteq [m] \) minimal for \( S \) if
    \[ \calK_{\calI} = \cbr{w : \tilde X_{\calI} w \succeq 0} \subseteq \calS \]
    and, for any \( j \in \calI \),
    \[
        \calK_{\calI \setminus j} \not \subseteq \calS.
    \]
    That is, removing any half-space constraint indexed by \( \calI \) ensures \( \calK_{\calI} \) is not contained in \( \calS \).
\end{definition}

Note that there may be many minimal index sets for a singular cone and these sets may have varying cardinalities.
However, each minimal index set shares a key property: every row \( \tilde x_i \) indexed by such \( \calI \) must be orthogonal to \( \calS \).

\begin{lemma}\label{lemma:orthogonality}
    Let \( \tilde X \in \R^{m \times d} \) such that the cone
    \( \calK = \cbr{ w: \tilde X w \succeq 0} \)
    is singular.
    Let \( \calS = \text{aff} (\calK) \) be the smallest containing subspace and \( \calI \) a minimal index set for \( \calS \).
    Then, \( \abr{\tilde x_i, s} = 0 \) for all \( i \in \calI \) and \( s \in \calS \).
\end{lemma}
\begin{proof}
    Suppose \( \abr{\tilde x_i, w} \neq 0 \) for some \( i \in \calI \) and \( w \in \calK \).
    Since \( w \in \calK \), \( \tilde X w \succeq 0 \) and it must be that \( \abr{\tilde x_i, w} > 0 \).
    Let \( z \in \calS^{\perp} \) be arbitrary and define \( w' = z + \alpha w \), \( \alpha > 0 \).
    By taking \( \alpha \) to be sufficiently large, we obtain
    \[
        \abr{\tilde x_i, w'} = \abr{\tilde x_i, z}  + \alpha \abr{\tilde x_i, w} > 0,
    \]
    Since \( w' \not \in \calS \supseteq \calK \), we must have
    \[ \tilde X_{\calI \setminus i} w' \not \succeq 0 \implies \tilde X_{\calI \setminus i} z \not \succeq 0, \]
    where we have used \( X w \succeq 0 \).
    Moreover, this holds for all \( z \in \calS^{\perp} \), which implies that \( \calK_{\calI \setminus i} \subseteq \calS \) and \( \calI \) cannot be minimal for \( \calS \).
    We conclude \( \abr{\tilde x_i, w} = 0 \) for all \( i \in \calI \) and \( w \in \calK \) by contradiction.

    Since \( \calS \) is the affine hull of \( \calK \), we have for every \( s \in \calS \) and \( i \in \calI \) the following:
    \[
        \abr{\tilde x_i, s} = \abr{\tilde x_i, \sum_{j=1}^k \alpha_j y_j} = \sum_{j=1}^k \alpha_j \abr{ \tilde x_i, y_j} = 0,
    \]
    since \( y_j \in \calK \).
\end{proof}

Similarly, if any constraint is tight at a relative interior point, then that constraint must be orthogonal to the cone.

\begin{lemma}\label{lemma:tight-constraints}
    Let \( \tilde X \in \R^{m \times d} \), \( \calK = \cbr{ w : \tilde X w \succeq 0} \), and \( \bar w \) be a relative interior point of \( \calK \).
    If \( \abr{\tilde x_j, \bar w} = 0 \) for any \( j \in [m] \), then \( \tilde x_j \) is orthogonal to \( \calK \).
\end{lemma}
\begin{proof}
    Suppose \( \abr{\tilde x_j, \bar w} = 0 \) for some \( j \in [m] \).
    If there exists \( w \in \calK \) such that \( \abr{\tilde x_j, \bar w} > 0 \), then \( \abr{\tilde x_j, \bar w + w} > 0 \) and \( \bar w + w \in \calK \), which contradicts the assumption \( \bar w \) is a relative interior point.
    Since every \( w \in \calK \) satisfies \( \abr{\tilde x_j, w} \geq 0 \), we conclude \( \abr{\tilde x_j, w} = 0 \) for all such \( w \).
\end{proof}

\cref{lemma:orthogonality} is key to our analysis because it implies that the half-space constraints which force \( \calK \) to lie in a subspace don't ``cut into'' that subspace.
In particular, it means that we can choose to enforce membership in \( \calH_{x_i} \) or \( \calH_{- x_i} \) without changing the inclusion.
We show now that there exists a choice of signed half-spaces for which the intersection is non-singular.

\begin{lemma}\label{lemma:sign-assignment}
    Let \( x_1, \ldots x_m \) be a collection of vectors in \( \R^d \) and \( X \in \R^{m \times d} \) the matrix formed by stacking these vectors.
    Then there exists a diagonal matrix \( \tilde D \), where \( \tilde D_{jj} \in \cbr{-1, 1} \), such that
    \[
        \text{aff} \rbr{\cbr{ w : \tilde D X w \succeq 0}} = \R^d.
    \]
\end{lemma}
\begin{proof}
    We proceed by induction.
    Let \( \tilde \calD_{11} = 1 \) and \( \calK_1 = \calH_{x_1} := \cbr{w : \abr{x_1, w} \geq 0} \).
    Clearly \( \text{aff} (\calK_1) = \R^d \) since it is a half-space.

    Now, let \( t < m \) and assume that \( \text{aff} (\calK_t) = \R^d \).
    Consider
    \[
        A_{t+1} := \calK_t \cap \calH_{x_{t+1}}.
    \]
    If \( \calS := \text{aff} (A_{t+1}) = \R^d \), then the inductive hypothesis holds at \( \calK_{t+1} = A_{t+1} \) and we can choose \( \tilde D_{t+1, t+1} = 1 \).
    Otherwise, \( x_{t+1} \) must be part of a minimal index set \( \calI \subseteq [t+1] \) such that \( \cbr{w : \tilde D_{\calI} X_{\calI} w \succeq 0 } \subseteq \calS \).
    \cref{lemma:orthogonality} now implies that \( x_{t+1} \) is orthogonal to \( \calS \).
    Let \( w \in \calK_{t} \cap \calS^{\perp} \) (which is non-empty by the inductive hypothesis) and observe that
    \[
        \abr{w, x_{t+1}} < 0,
    \]
    must hold, otherwise \( w \in A_{t+1} \).
    We deduce \( \abr{w, x_{t+1}} \leq 0 \) for every \( w \in \calK_{t} \) and thus
    \[
        \calK_t \cap \calH_{- x_{t+1}} = \calK_t,
    \]
    which is full-dimensional by the inductive hypothesis.
    Taking \( \calK_{t+1} = \calK_t \) and \( \calD_{t+1, t+1} = -1 \) completes the case.

    The desired result follows by induction.
\end{proof}

We now use \cref{lemma:sign-assignment} to show that every singular cone is contained in a non-singular cone.

\coneContainment*
\begin{proof}
    For simplicity, we drop the index \( i \) and work with \( \calK = \cbr{w : (2 D - I) X w \succeq 0 } \).
    To ease the notation, we also write \( \tilde D = (2 D - I) \) and \( \tilde X = \tilde D X \).
    Let \( \calS = \text{aff} (\calK) \) be the smallest subspace containing \( \calK \),
    \( \calO \) be the set of all indices \( j \) such that \( x_j \) is orthogonal to \( \calS \),
    and \( \calU = [n] \setminus \calO \).
    \cref{lemma:sign-assignment} implies that there exists an alternative activation pattern \( \tilde D' \) such that,
    \( \tilde D_{\calU}' = \tilde D_{\calU} \), and
    if \( \calK' = \cbr{w : \tilde D' X w \succeq 0 } \), then
    \[
        \calK'_{\calO} := \cbr{w : \tilde D_{\calO}' X_\calO w \succeq 0} \quad \text{satisfies} \quad \text{aff} \rbr{\calK_\calO'} = \R^d.
    \]
    Since the vectors indexed by \( \calO \) are orthogonal to \( \calS \), they are also orthogonal to every \( w \in \calK \), implying \( \calK \subset \calK' \).
    In other words, the change of activation signs preserves inclusion of \( \calK \).

    Let us show that \( \calK' \) contains an interior point.
    Let \( \bar w \) be a relative interior point of \( \calK \) and suppose that \( \abr{\tilde x_j, \bar w} = 0 \) for some \( j \in \calU \).
    \cref{lemma:tight-constraints} implies \( \tilde x_j \) is orthogonal to \( \calK \).
    But, then \( j \in \calO \), which is a contradiction.
    We conclude \( \tilde X_{\calU }\bar w \succ 0\).

    Let \( \bar y \) be an interior point of \( \calK'_{\calO} \), which exists because \( \text{aff} (\calK'_{\calO}) = \R^d \).
    The point \( \bar z = \bar y + \alpha \bar w \), \( \alpha > 0 \) satisfies
    \[
        \tilde D_{\calO}' X_{\calO} \bar z = \tilde D_{\calO}' X_{\calO} \bar y \succ 0,
    \]
    since \( \bar w \) is orthogonal to the rows of \( X_{\calO} \).
    Similarly, by taking \( \alpha \) to be sufficiently large, we have
    \[
        \tilde D_{\calU}' X_{\calU} \bar z \succ 0.
    \]
    We have shown that \( \bar z \) is an interior point of \( \calK' \)
    and \cref{lemma:affine-characterization} now implies \( \calK' - \calK' = \R^d \).
\end{proof}

\coneElimination*
\begin{proof}
    Assume \( \lambda = 0 \).
    The equivalence of the C-ReLU and C-GReLU problems with sub-sampled
    patterns \( \tilde \calD \setminus \calS(\tilde\calD) \) follows
    immediately from cone decompositions.
    Let \( w^* \) be a solution to the C-GReLU problem with
    \( \tilde \calD \setminus \calS(\tilde\calD) \)
    and \( p^* \) the optimal value.
    Let \( d^* \) be the optimal value of the C-ReLU problem with the same
    patterns.
    For every \( D_i \in \tilde \calD \setminus \calS(\tilde\calD) \),
    \( \calK_i - \calK_i = \R^d \) by construction.
    Thus, there exists \( v_i, u_i \in \calK_i \) such that
    \( v_i - u_i = w_i^* \).
    It is now straightforward to conclude that,
    \[
        p^* =
        \calL\rbr{\sum_{D_i \in \tilde \calD \setminus \calS(\tilde\calD)}
            D_i X w_i^*, y}
        =
        \calL\rbr{\sum_{D_i \in \tilde \calD \setminus \calS(\tilde\calD)}
            D_i X (v_i - u_i), y}
        \geq d^*,
    \]
    since \( v_i, u_i \) are feasible for the C-ReLU problem.
    Letting \( v^*, u^* \) be a solution to the C-ReLU problem,
    we obtain the opposite inequality immediately be noting that
    \( w_i = v^*_i - u^*_i \) is feasible for the C-GReLU problem.
    Thus, the two problems are equivalent when singular cones are omitted.
    Now let us show that singular cones do not affect the optimal value
    of the C-ReLU problem.

    Assume \( \lambda \geq 0 \) and let \( (v^*, u^*) \) be a solution to the
    C-ReLU problem with patterns \( \tilde\calD \) and \( p^* \) the optimal
    value.
    For every \( D_i \in \calD_{X} \), we either have \( \calK_i - \calK_i = \R^d \), or not.
    If the former condition holds, take \( r_i^* = v_i^* \) and \( q_i^* = u_i^* \).
    Otherwise, invoke \cref{prop:cone-containment} to obtain \( D_j \in \calD_{X} \)
    such that \( \calK_i \subset \calK_j \) and \( \calK_j - \calK_j = \R^d \).
    By the construction in the proof of \cref{prop:cone-containment},
    if \( \sbr{D_j}_{kk} \neq \sbr{D_i}_{kk} \),
    then \( x_k \) is orthogonal to \( v_i^*, u_i^* \in \calK_i \) and we deduce
    \[
        D_j X (v_i^* - u_i^*)  = D_i X (v_i^* - u_i^*).
    \]
    We may therefore merge the two neurons as \( r^*_j = v^*_j + v^*_i \),
    \( q_j^* = w_j^* + u_i^* \) and update
    \( \tilde \calD' = \tilde \calD \setminus D_i \) without changing the
    (optimal) prediction of the C-ReLU model.
    In this way, we obtain a sub-sampled C-ReLU problem with activation patterns
    \( \tilde \calD \setminus \calS(\tilde \calD) \),
    for which
    \begin{align*}
        p^* & = \calL\rbr{\sum_{D_i \in \tilde \calD} D_i X (v_i^* - u_i^*), y}
        + \lambda \sum_{D_i \in \tilde \calD} \norm{v_i^*}_2 + \norm{v_i^*}_2
        \\
            & \geq \calL\rbr{\sum_{D_j \in \tilde \calD \setminus \calS(\tilde \calD)} D_j X (r_j^* - q_j^*), y}
        + \lambda \sum_{D_j \in \tilde \calD \setminus \calS(\tilde \calD)} \norm{r_j^*}_2 + \norm{q_j^*}_2                              \\
            & \geq \min_{r_j, q_j \in \calK_j} \calL\rbr{\sum_{D_j \in \tilde \calD \setminus \calS(\tilde \calD)} D_j X (r_j - q_j), y}
        + \lambda \sum_{D_j \in \tilde \calD \setminus \calS(\tilde \calD)} \norm{r_j}_2 + \norm{q_j}_2
        := d^*,
    \end{align*}
    where we used triangle inequality to imply \( \norm{r_j^*}_2 \leq
    \norm{v_j^*}_2 + \norm{v_i^*}_2 \) also the fact that \( r_i^*, q_i^* \)
    are feasible because \( \calK_j \) is a convex cone and closed under
    addition.
    It is now straightforward to conclude \( p^* = d^* \) since sub-sampling
    the C-ReLU problem can only increase the optimal objective value.
\end{proof}

\subsection{Approximating ReLU by Cone Decomposition: Proofs}\label{app:cone-decompositions}
\closedFormDecomp*
\begin{proof}
    Let's show that the decomposition is valid.
    Setting \( v = u + w \),
    we obtain \( v - w = u \) by construction.
    For notational ease, let \( \calJ = [n] \setminus \calI \).
    It holds that
    \[
        \tilde X_{\calI} w = - \tilde X_{\calI} \tilde X_{\calI}^\dagger \tilde X_{\calI} u = - \tilde X_{\calI} u > 0.
    \]
    Moreover, since \( X \) is full row-rank and \( \calJ \cap \calI = \emptyset \), we have
    \[
        \tilde X \tilde X^\dagger = \tilde X \tilde X^\top \rbr{\tilde X \tilde X^\top}^{-1} = I,
    \]
    which implies that \( \tilde X_{\calJ} \tilde X_{\calI}^\dagger = 0 \).
    We deduce
    \[
        \tilde X_{\calJ} w = - \tilde X_{\calJ} \tilde X_{\calI}^\dagger \tilde X_{\calI} u = 0.
    \]
    Moreover, we also have
    \[
        \tilde X_{\calI} v = \tilde X_{\calI} u - \tilde X_{\calI} u = 0,
    \]
    and
    \[
        \tilde X_{\calJ} v = \tilde X_{\calJ} u + 0 \geq 0,
    \]
    by definition of \( \calJ \).
    We conclude \( w, v \in \calK \).

    To show the approximation result, start from
    \begin{align*}
        \norm{w}_2
         & = \norm{\tilde X_{\calI}^\dagger \tilde X_{\calI} u}_2             \\
         & \leq \norm{\tilde X_{\calI}^\dagger \tilde X_{\calI}}_2 \norm{u}_2 \\
         & = \norm{u}_2.
    \end{align*}
    Triangle inequality now gives,
    \begin{align*}
        \norm{v}_2
         & = \norm{u - \tilde X_{\calI}^\dagger \tilde X_{\calI} u}_2               \\
         & \leq \norm{(I - \tilde X_{\calI}^\dagger \tilde X_{\calI})}_2 \norm{u}_2 \\
         & = \norm{u}_2,
    \end{align*}
    and summing these two inequalities gives the result.
\end{proof}

\worstCaseNorm*
\begin{proof}
    Consider the data matrix
    \begin{equation}
        X =
        \begin{bmatrix}
            1   & \alpha \\
            - 1 & \alpha \\
        \end{bmatrix}
    \end{equation}
    The cone corresponding to positive activations for both examples is
    \[
        \calK = \cbr{x \in \R^d : x_2 \geq - x_1 / \alpha, x_2 \geq x_1 / \alpha  }.
    \]
    Consider decomposing the vector \( w = [2, 0] \) onto \( \calK - \calK \).
    Clearly \( w \not \in \calK \);
    by inspection, we see that the minimum norm decomposition is given by
    \( v = [1, 1/\alpha] \), and \( u = [-1, 1/\alpha] \).
    Taking \( \alpha \rightarrow 0 \), we find \( \norm{v}_2 = \norm{u}_2 \rightarrow \infty \).
\end{proof}

\socpDecomp*
\begin{proof}
    First, we re-parameterize the problem: \( w = v - u \) implies \( u + w = v \),
    giving the equivalent program
    \begin{equation}
        \begin{aligned}
            \min_{u \in \calK} \norm{u}_2 + \norm{u + w}_2^2.
        \end{aligned}
    \end{equation}
    In order to characterize the solution, we re-write the constraints into a
    single system of linear inequalities as follows:
    \begin{align*}
        \calF
              & = \cbr{u : \tilde X u \succeq 0, \tilde X u \succeq - \tilde X w} \\
              & = \cbr{u : \tilde X u \succeq 0, \tilde X u \succeq b},
        \intertext{where we have introduced \( b = - \tilde X w \). It is possible to combine these inequalities by taking the element-wise maximum as follows:}
        \calF & = \cbr{u : \tilde X u \succeq \max\rbr{0, b}}                     \\
              & = \cbr{u : \tilde X u \succeq \rbr{b}_+}.
    \end{align*}
    Let \( \bar u \in \calF \) be a optimal point for the reparameterized program.
    Relaxing the objective using triangle inequality gives,
    \begin{align*}
        \norm{\bar u}_2 + \norm{w + \bar u}_2
         & = \min_{u \in \calF} \norm{u}_2 + \norm{u + w}_2  \\
         & \leq \min_{u \in \calF} 2\norm{u}_2 + \norm{w}_2.
    \end{align*}
    Let \( u' \) be a solution to the relaxation.
    The KKT conditions imply there exists a submatrix \( \tilde X_{\calJ} \)
    for which the inequality constraints are tight:
    \[
        \tilde X_{\calJ} u' = \rbr{b_{\calJ}}_+.
    \]
    The set of vectors satisfying this equality is 
    \( \cbr{\sbr{\tilde X_{\calJ}}^\dagger \rbr{b_{\calJ}}_+ + z : z \in \text{null} (X_{\calJ})} \).
    Choosing \( z \neq 0 \) can only increase the value of the objective, from which we deduce \( w' = \sbr{\tilde X_{\calJ}}^\dagger \rbr{b_{\calJ}}_+ \).
    We obtain
    \begin{align*}
        \norm{\bar u}_2 + \norm{\bar v}_2
         & \leq 2\norm{u'}_2 + \norm{w}_2                                                                                  \\
         & = 2 \norm{\sbr{\tilde X_{\calJ}}^\dagger \rbr{b_{\calJ}}_+} + \norm{w}_2                                        \\
         & \leq \frac{2}{\sigma_{\text{min}}(\tilde X_{\calJ})} \norm{\rbr{b_{\calJ}}_+}_2 + \norm{w}_2                    \\
         & \leq \frac{2}{\sigma_{\text{min}}(\tilde X_{\calJ})} \norm{b_{\calJ}}_2 + \norm{w}_2                            \\
         & =  \frac{2}{\sigma_{\text{min}}(\tilde X_{\calJ})} \norm{\tilde X_{\calJ} w}_2 + \norm{w}_2                     \\
         & \leq \rbr{1 + 2\frac{\sigma_{\text{max}}(\tilde X_{\calJ})}{\sigma_{\text{min}}(\tilde X_{\calJ})}} \norm{w}_2.
    \end{align*}
\end{proof}

\approxResult*
\begin{proof}
    The proof is straightforward given our existing results.
    Let \( w^* \) be the solution to the full (potentially regularized) C-GReLU problem.
    For every \( D_i \in \calD_{X} \), we either have \( \calK_i - \calK_i = \R^d \), or not.
    If the former condition holds, take \( r_i^* = w_i^* \).
    Otherwise, invoke \cref{prop:cone-containment} to obtain \( D_j \in \calD_{X} \) such that \( \calK_i \subset \calK_j \) and \( \calK_j - \calK_j = \R^d \).
    By the construction in the proof of \cref{prop:cone-containment}, if \( \sbr{D_j}_{kk} \neq \sbr{D_i}_{kk} \), then \( x_k \) is orthogonal to \( \calK_i \) and we deduce
    \[
        D_j X w^*_i = D_i X w^*_i.
    \]
    We may therefore merge the two neurons as \( r^*_j = w^*_j + w^*_i \) and update 
    \( \tilde \calD' = \tilde \calD \setminus D_i \) without changing the 
    loss component of the C-GReLU program.
    Furthermore, since \( \norm{r^*_j} \leq \norm{w^*_j}_2 + \norm{w^*_i}_2 \), 
    merging these neurons can only decrease the regularization term.
    \footnote{In fact, we know that one of \( w^*_j, w^*_i \) is zero or they are collinear}
    In this way, we obtain a sub-sampled C-GReLU problem with activation
    patterns \( \tilde \calD \) and optimal value \( d^* \).

    Let \( \rbr{v^*, u^*} \) be the optimal solution to full C-ReLU problem.
    Applying \cref{prop:socp-decomp} for each \( D_i \in \tilde \calD \) gives decompositions \( w_i^* = v_i' - u_i' \) such that
    \begin{align*}
        d^*
         & = \mathcal{L}\Big(  \sum_{D_i \in \tilde \calD} D_i X w_i^*, y\Big) + \lambda \sum_{D_i \in \tilde \calD} \norm{w_i^*}_2                                                                                                                                               \\
         & \leq \mathcal{L}\Big(  \sum_{D_i \in \tilde \calD} D_i X v^*_i - u_i^*, y\Big) + \lambda \sum_{D_i \in \tilde \calD} \norm{v_i^* - u_i^*}_2                                                                                                                            \\
         & \leq \mathcal{L}\Big(  \sum_{D_i \in \tilde \calD} D_i X v^*_i - u_i^*, y\Big) + \lambda \sum_{D_i \in \tilde \calD} \norm{v_i^*}_2 + \norm{u_i^*}_2                                                                                                                   \\
         & = p^*                                                                                                                                                                                                                                                                  \\
         & \leq \mathcal{L}\Big(  \sum_{D_i \in \tilde \calD} D_i X v'_i - u_i', y\Big) + \lambda \sum_{D_i \in \tilde \calD} \norm{v_i'}_2 + \norm{u_i'}_2                                                                                                                       \\
         & = \mathcal{L}\Big(  \sum_{D_i \in \tilde \calD} D_i X w^*_i, y\Big) + \lambda \sum_{D_i \in \tilde \calD} \norm{v_i'}_2 + \norm{u_i'}_2                                                                                                                                \\
         & \leq \mathcal{L}\Big(  \sum_{D_i \in \tilde \calD} D_i X w^*_i, y\Big) + \lambda \sum_{D_i \in \tilde \calD} \norm{w_i^*}_2 + \lambda \sum_{D_i \in \tilde \calD} 2 \frac{\sigma_{\text{max}}(\tilde X_{\calJ})}{\sigma_{\text{min}}(\tilde X_{\calJ})} \norm{w_i^*}_2 \\
         & = p^* + \lambda \sum_{D_i \in \tilde \calD} 2 \frac{\sigma_{\text{max}}(\tilde X_{\calJ})}{\sigma_{\text{min}}(\tilde X_{\calJ})} \norm{w_i^*}_2.
    \end{align*}
    We have abused notation here and omitted the dependence on \( i \) in \( \tilde X_{\calJ} = \rbr{2 D_i -I} X_{\calJ_i} \).
    However, observe that \( 2D_i - I \) is orthonormal so that \( \sigma_{\text{max}}(\tilde X_{\calI}) = \sigma_{\text{max}}(X_{\calI}) \) and \( \sigma_{\text{min}}(\tilde X_{\calI}) = \sigma_{\text{min}}(X_{\calI}) \) for all \( \calI \).
    Maximizing over \( \calI \subseteq [n] \) now gives a fixed subset \( \calJ \) for which the claimed bound holds.
\end{proof}

\iffalse

    %% Now superseded by other results!
    \subsection{Additional Results on Cone Decompositions}\label{app:additional-cone-results}

    In this section, we provide additional theoretical results which were omitted from the main paper due to space constraints.
    We examine the special case where \( X \) is full-row rank and also \emph{semi-orthogonal}, meaning the rows of \( X \) are orthogonal (although not necessarily orthonormal) vectors.
    In this setting, there exists a simple decomposition based on the pseudo-inverse of a submatrix of \( X \) that gives a \( 2 \)-approximation of the norm.

    \begin{proposition}\label{prop:orthogonal-decomp}
        Let \( X \) be full row-rank and semi-orthogonal.
        For every \( u \in \R^d \),
        if \( \calI = \cbr{i \in [n] : \abr{\tilde x_i, u} < 0} \), then the decomposition
        \[
            w = - \tilde X_{\calI}^\dagger \tilde X_{\calI} u,
        \]
        and \( v = u + w \) satisfies \( v, w \in \calK \) and
        \[
            \norm{v}_2 + \norm{w}_2 \leq 2 \norm{u}_2.
        \]
    \end{proposition}
    \begin{proof}
        Let's show that the decomposition is valid.
        Clearly, \( v - w = u \) by construction.
        For notational ease, let \( \calJ = [n] \setminus \calI \).
        It holds that
        \[
            \tilde X_{\calI} w = - \tilde X_{\calI} X_{\calI}^\dagger \tilde X_{\calI} u = - \tilde X_{\calI} u > 0,
        \]
        and, using that fact that \( \tilde X_{\calI}^\dagger X_{\calI} u_i^* \) lies in the row space of \( \tilde X_{\calI} \), which is orthogonal to the row space of \( \tilde X_{\calJ} \),
        \[
            \tilde X_{\calJ} w = - \tilde X_{\calJ} \tilde X_{\calI}^\dagger \tilde X_{\calI} u = 0.
        \]
        It remains only to note that
        \[
            \tilde X_{\calI} v = \tilde X_{\calI} u - \tilde X_{\calI} u = 0,
        \]
        and
        \[
            \tilde X_{\calJ} v = \tilde X_{\calJ} u + 0 \geq 0,
        \]
        by definition of \( \calJ \).
        We conclude that \( w, v \in \calK \).

        To show the approximation result, start from
        \begin{align*}
            \norm{w}_2
             & = \norm{\tilde X_{\calI}^\dagger \tilde X_{\calI} u}_2             \\
             & \leq \norm{\tilde X_{\calI}^\dagger \tilde X_{\calI}}_2 \norm{u}_2 \\
             & = \norm{u}_2.
        \end{align*}
        Triangle inequality now gives the final result as follows:
        \begin{align*}
            \norm{v}_2
             & = \norm{u - \tilde X_{\calI}^\dagger \tilde X_{\calI} u}_2               \\
             & \leq \norm{(I - \tilde X_{\calI}^\dagger \tilde X_{\calI})}_2 \norm{u}_2 \\
             & = \norm{u}_2.
        \end{align*}
        Combining these two inequalities gives the result.
    \end{proof}
\fi

%% Old Proof

%\begin{proof}
%    We first show that the decomposition is valid.
%    First, we observe that \( Z = \tilde X (\tilde X \tilde X^\top)^{-1} \) so that
%    \[
%        \tilde X Z^\top = I,
%    \]
%    and \( \tilde X_{\calI} \tilde Z_{\calI}^{\top} \) is diagonal with diagonal values
%    exactly equal to \( 1 \).
%    Using the index set \( \calI \), we have
%    \begin{align*}
%        \tilde X_{\calI} w
%         & = - \tilde X_{\calI} \tilde Z_{\calI}^{\top} \rbr{\tilde X_{\calI} \tilde Z_{\calI}^{\top}}^{-1} \tilde X_{\calI} u \\
%         & = - \tilde X_{\calI} u                                                                                              \\
%         & > 0,
%    \end{align*}
%    by choice of \( \calI \).
%    Moreover, since \( \tilde X_{[n] \setminus \calI} \tilde Z_{\calI}^{\top} = 0 \),
%    we have \( \tilde X_{[n] \setminus \calI} w = 0 \) and \( w \in \calK \).
%    The other component similarly satisfies
%    \begin{align*}
%        \tilde X_{[n] \setminus \calI} v
%         & = \tilde X_{\calI} u - \tilde X_{\calI} u \\
%         & = 0
%        \intertext{and}
%        \tilde X_{[n]\setminus\calI} v
%         & = \tilde X_{[n]\setminus\calI} u - 0      \\
%         & \geq 0,
%    \end{align*}
%    again by choice of \( \calI \).
%    Thus, \( v \in \calK \) and the decomposition is valid.
%    Next we control the blow-up in norm caused by the decomposition.

%    Let \( \sigma_{\text{min}} \) be the smallest non-zero singular value of \( \tilde X_{\calI} \) and \( \sigma_{\text{max}} \) the largest singular value.
%    Observe that \( \abs{\abr{\tilde x_j, \tilde z_j}} = \norm{\tilde X_{\calI} \tilde z_j}_2 \geq \sigma_{\text{min}} \) for \( j \in \calI \) since \( \abr{\tilde x_l, \tilde z_j} = 0 \) for \( l \neq j \) and \( \norm{\tilde z_j}_2 = 1 \).
%    Using orthogonality of the \( \tilde z_j \), we obtain
%    \begin{align*}
%        \norm{w}_2^2
%                   & = \left\|\tilde Z_{\calI}^{\top} \rbr{\tilde X_{\calI} \tilde Z_{\calI}^{\top}}^{-1} \tilde X_{\calI} u\right\|_2^2                                          \\
%                   & = \left\|\sum_{j \in \calI} \frac{\abr{\tilde x_j, u}^2}{\abr{\tilde x_j, \tilde z_j}^2} \tilde z_j \right\|_2^2                                             \\
%                   & = \sum_{j \in \calI} \frac{\abr{\tilde x_j, u}^2}{\abr{\tilde x_j, \tilde z_j}^2}                                \tag{since \( \norm{\tilde z_j}_2^2 = 1 \)} \\
%                   & \leq \frac{1}{\sigma_{\text{min}}^2}  \sum_{j \in \calI} \abr{\tilde x_j, u}^2                                                                               \\
%                   & = \frac{1}{\sigma_{\text{min}}^2}  \norm{\tilde X_{\calI} u}_2^2                                                                                             \\
%                   & \leq \frac{\sigma_{\text{max}}^2}{\sigma_{\text{min}}^2}  \norm{u}_2^2.                                                                                      \\
%        \implies
%        \norm{w}_2 & \leq \frac{\sigma_{\text{max}}}{\sigma_{\text{min}}}  \norm{u}_2.
%    \end{align*}
%    Similarly,
%    \begin{align*}
%        \norm{v}_2
%         & = \norm{u + w}_2                                                             \\
%         & \leq \norm{u}_2 + \norm{w}_2                                                 \\
%         & \leq \rbr{1 + 2 \frac{\sigma_{\text{max}}}{\sigma_{\text{min}}}} \norm{u}_2.
%    \end{align*}
%    Combining these two equations gives the desired result.
%\end{proof}

